\chapter{Feasibility Report}
\label{ch:feasibility}

A feasibility study asks whether a system can be built and used successfully
under real technical, operational, economic, schedule and legal
constraints~\cite{sommerville2016,pressman2015}. At proposal stage, feasibility
can only be argued. At the mid-term point, much of FoodLink exists, so this
chapter assesses feasibility against evidence from the implementation. It
distinguishes throughout between three categories:
\begin{itemize}
  \item what is \textbf{already implemented} and verified;
  \item what is \textbf{feasible but not implemented}, meaning the design
    accommodates it and no technical obstacle is known;
  \item what remains \textbf{future work} requiring further design or
    decisions.
\end{itemize}

\section{Technical Feasibility}
\label{sec:feas-technical}

\subsection{Technology Stack}
Every component of the stack is mature, widely used and well documented
(Table~\ref{tab:stack}): Python with FastAPI, SQLAlchemy, Alembic and
Pydantic on the server; React, TypeScript and Vite in the browser; pytest and
Vitest for testing~\cite{fastapiDocs,sqlalchemyDocs,alembicDocs,pydanticDocs,reactDocs,viteDocs,pytestDocs,vitestDocs}.
The project needed no experimental technology, no proprietary component and
no external service. The strongest evidence of technical feasibility is that
the complete workflow from posting to completion runs today and passes
651 automated tests (Section~\ref{sec:metrics}).

\subsection{Frontend}
\emph{Implemented.} Four role portals are built with standard React patterns:
context providers, a router and a single HTTP client. The browser performs
only presentation work, plus one computation that could not live on the server,
namely resizing a photograph that only the browser holds. Frontend
feasibility risks are therefore low.
\emph{Feasible but not implemented.} Routing phone visitors to the existing
\code{/m/*} layouts automatically (a detection hook exists but is unused),
generating the TypeScript API types from the server's OpenAPI document instead
of maintaining them by hand, and adding controls for the API-only operations.
\emph{Future work.} Code-splitting per portal and a caching data layer, if
profiling shows the full-list reloads to be too slow.

\subsection{Backend and APIs}
\emph{Implemented.} 27 REST routes cover every proposal workflow. The API is
stateless apart from rate-limit counters, documents itself through OpenAPI,
and expresses its rules as inspectable data: transition, role and scope
tables. Complexity was concentrated where it is justified, in authorization
scopes, conditional writes and disclosure rules, rather than in
infrastructure.
\emph{Feasible but not implemented.} A global exception handler that returns a
readable message and a correlation id; bounds on profile text fields; a
scheduled caller for the expiry sweep. The sweep is already an idempotent
endpoint, so a cron job holding an administrator token would suffice.

\subsection{Database}
\emph{Implemented.} Six normalised tables under Alembic migration control.
Timestamps are normalised to UTC. SQLite needs no database server, which kept
development setup to installing packages and setting one environment variable.
\emph{Feasible but not implemented.} Moving to PostgreSQL. All data access
goes through SQLAlchemy, and the lifecycle's concurrency guards are plain
conditional \code{UPDATE} statements that need no engine-specific feature.
However, no PostgreSQL driver is pinned, the suite has never run against
PostgreSQL, and the guards' behaviour there depends on PostgreSQL's default
\textsc{read committed} isolation level (D-28). Enabling foreign-key
enforcement in SQLite is a small, localised change.
\emph{Constraint.} SQLite permits one writer at a time~\cite{sqliteWhenToUse}.
It is suitable for a single-process pilot but not for several API worker
processes.

\subsection{Authentication}
\emph{Implemented.} Password hashing with bcrypt, signed expiring tokens, a
fail-closed signing key, per-request re-reading of the account, and
administrator-only creation of administrators.
\emph{Feasible but not implemented.} Token revocation (for example a version
number on the user row checked on each request), shorter-lived access tokens
with refresh, and security headers.
\emph{Future work.} Email verification, multi-factor authentication and an
administrator audit log, which need product decisions and, for email, an
external service.

\subsection{Deployment Considerations}
\emph{Not implemented:} there is no deployment configuration of any kind.
Hosting FoodLink is technically straightforward, since it is one ASGI
application and one static bundle. The implementation, however, fixes several
decisions that a deployment must respect:
\begin{enumerate}
  \item migrations currently run inside the application's start-up, which is
    safe only while exactly one process runs (D-23);
  \item rate-limit counters live in process memory, so several processes would
    multiply the effective limits unless a shared store is introduced (D-27);
  \item behind a reverse proxy, the server must be configured to trust the
    proxy's forwarded client address, or all users will share one rate-limit
    budget;
  \item \code{CORS_ORIGINS} must list the deployed frontend's origin, and a
    real signing key must be supplied as a secret;
  \item the expiry sweep needs an external scheduler.
\end{enumerate}
A single-process deployment with SQLite is feasible now. A multi-process
deployment is feasible once items 1 and 2 are addressed and PostgreSQL is
validated.

\subsection{Implementation Complexity}
The domain is small (six tables, nine states), but the implementation is not
trivial. Much of the effort went into properties that tutorials rarely cover:
per-role read scopes expressed as SQL, disclosure rules that change what a
record says to different readers, race-safe transitions, and privacy
protections against location inference. The project has shown that a
four-person student team can deliver this complexity by working in small,
tested increments with written decisions. The remaining work
(Section~\ref{sec:feas-schedule}) is of lower algorithmic complexity than what
has already been built.

\textbf{Assessment:} technically feasible, and already demonstrated for the
core system. The main open technical risk is an untested production database
and deployment topology.

\section{Operational Feasibility}
\label{sec:feas-operational}

Operational feasibility asks whether the intended users can and will use the
system as designed.

\paragraph{Donors.} A donor signs up, opens \emph{Create Food Donation}, fills in
one form and submits it. The photo is optional and resized automatically. The
deadline is entered as a clock time, and a time that has already passed today
is taken to mean tomorrow. The pickup pin comes from the browser's location
service or is typed as latitude and longitude; there is no map picker, which
is a usability gap for non-technical donors. The donor can then follow the
donation's timeline and browse the needs board to see what verified
organisations want. The flow suits a kitchen worker posting during clean-up,
as the proposal required. Two operational gaps remain. A donor must re-open
the portal to see progress, since nothing is pushed. A donor who needs to
withdraw a donation has no button for it yet.

\paragraph{Recipient organisations.} An organisation registers and then waits
for an administrator's verification. Until verified, it sees an explanation
rather than an empty screen. Once verified, its \emph{Available Donations}
screen lists collectable donations, each with its own score and plain-language
reasons, and acceptance is a single action. It confirms receipt from
\emph{Accepted Deliveries} and maintains standing needs on
\emph{Food Requirements \& Demand Signals}. Verification adds a deliberate
human delay before an
organisation becomes useful. That is the cost of the trust model, and an
administrator must actually carry out the vetting. A more serious operational
gap is that no screen lets an organisation set its map coordinates. A newly
registered organisation can browse and accept once verified, but it is
never \emph{suggested} to donors until its pin is set through the API.
Adding coordinate fields (or a location button like the donor's) to the
organisation profile is a small, high-priority change.

\paragraph{Volunteer couriers.} A courier sees unclaimed pickups with a coarse
area, enough to judge whether a run is practical, and claims one. The exact
address and donor appear only after the claim. The courier then records pickup
and delivery with one button each, and a history page lists completed work.
The availability toggle is advisory: it is recorded but does not currently
gate claiming. If a courier cannot complete a claimed run, only the accepting
organisation (through the API) or an administrator can release it at present.

\paragraph{Administrators.} Administration involves verifying organisations,
monitoring donations, running the expiry sweep and reviewing analytics. These
tasks are supported by portal screens. Account suspension and administrator
creation are available only through the API or the command-line tool, which is
workable for a small pilot run by technical staff but not for a larger service.

\paragraph{Practicality of the workflow.} The workflow keeps a human decision
at every point of custody. An organisation must accept, a courier must claim,
and the organisation must confirm receipt. The system never assigns food
automatically, which matches how the proposal expected real handovers to
work~\cite{foodlinkProposal}. The main operational risk is \emph{timeliness
without notifications}. The platform removes the need to find a recipient, but
it cannot yet tell that recipient that food is waiting.

\textbf{Assessment:} operationally feasible for a supervised pilot on one
campus. Notifications, portal controls for cancellation and release, and an
account-management screen are needed before unsupervised use.

\section{Economic Feasibility}
\label{sec:feas-economic}

This section identifies cost drivers without inventing monetary figures,
since no hosting provider, pilot size or service contract has been chosen.

\paragraph{Development resources.} Development effort is the time of the
four-member student team within the course. No paid developers, contractors
or consultants are involved.

\paragraph{Software and tooling.} Every runtime dependency and development
tool in the repository is open-source software, available without licence
fees. The only services used are a Git hosting platform (for the repository,
CI runs and the course documentation pages) and an online LaTeX editor for
reports. The report's cost assumptions do not
depend on any paid tier of either.

\paragraph{Infrastructure.} There is currently no infrastructure cost, because
nothing is deployed. The backend makes no calls to paid external APIs (no
maps, SMS or email providers), so operating costs would be limited to
hosting. A pilot would need, at minimum, one small application server and,
later, a managed PostgreSQL database and static hosting for the frontend.
Features on the roadmap would add costs. Email or SMS notifications carry
per-message charges; object storage for photos and road-distance routing
typically carry usage-based charges. Their magnitude depends on provider and
volume and has not been estimated.

\paragraph{Cost-reducing design choices already made.} SQLite by default
avoids a database server during development and a small pilot. Straight-line
distance avoids a routing API. Browser-side photo resizing reduces storage and
bandwidth. No paid third-party service is integrated.

\paragraph{Benefit side.} The proposal's intended benefit is surplus food
reaching people instead of being discarded~\cite{foodlinkProposal}. The
platform can now measure this through rescue rate and expiry-loss rate once it
is used. No benefit figures are claimed here, because none have been measured.

\textbf{Assessment:} economically feasible at course and pilot scale, with
near-zero software cost. Hosting and messaging costs for a wider service
remain to be estimated once a provider and scale are chosen.

\section{Schedule Feasibility}
\label{sec:feas-schedule}

\paragraph{Course timeline.} The course schedules a proposal in week~4, this
prototype-stage report in week~7, and the final report and seminar in
week~17~\cite{ucs503Project}.

\paragraph{Progress so far.} The repository's first commit is dated 29 July
2026. By commit \texttt{e588b35} (13 September 2026) it contained 85 commits.
The backend and frontend were built, mock data was removed, and two audit
cycles were run and their P1 findings fixed. Table~\ref{tab:schedule}
summarises implementation maturity by area.

\begin{table}[htbp]
\centering
\small
\begin{tabular}{|L{4.0cm}|L{3.2cm}|L{6.1cm}|}
  \hline
  \textbf{Area} & \textbf{Maturity} & \textbf{Remaining} \\
  \hline
  Donation lifecycle & Complete, hardened & Guard the expiry sweep; admin expiry accounting \\
  \hline
  Matching and explainability & Complete & Tuning weights needs outcome data \\
  \hline
  Authentication and authorization & Complete for pilot & Revocation, security headers, profile bounds \\
  \hline
  Requirements and needs board & Complete & Structured food categories (optional) \\
  \hline
  Web portals & Complete for core flows & Cancel, release and account-management controls; deep-link fix \\
  \hline
  Phone layouts & Partial & Remove demonstration residue; entry routing \\
  \hline
  Testing and CI & Strong backend; partial frontend & Frontend tests for data context, actions and forms \\
  \hline
  Deployment & Not started & Configuration, PostgreSQL validation, scheduler, secrets \\
  \hline
  Pilot and measurement & Not started & Recruit users; collect ledger metrics \\
  \hline
\end{tabular}
\caption{Implementation maturity and remaining work by area}
\label{tab:schedule}
\end{table}

\paragraph{Remaining work before the final evaluation.} In priority order, as
recorded in the project task list: the remaining P2 items (a global
exception handler, profile field bounds, the deep-link bug, frontend test
gaps and mobile residue); portal controls for the API-only operations; a
deployment configuration with a scheduled expiry sweep; and, if deployment
succeeds, a short pilot to produce real values for the proposal's metrics.
Each item is small or medium in size compared with the work already
completed.

\paragraph{Realistic future development areas.} Beyond the course timeline:
notifications, road distances, object storage for photos, weight tuning from
outcome data, and broader organisation onboarding.

\textbf{Assessment:} schedule-feasible. The core system was delivered by
mid-term. The remaining course-scope work is bounded. The principal schedule
risk is the deployment and pilot, which depend on arranging real participants
as well as on engineering.

\section{Legal, Security and Privacy Feasibility}
\label{sec:feas-legal}

FoodLink stores personal data about identifiable people: names, email
addresses, phone numbers, organisation addresses, donors' pickup locations
(which may be homes), photographs and activity history. Any real deployment in
India would have to comply with the Digital Personal Data Protection Act,
2023~\cite{dpdpAct2023}. This report does not offer a legal opinion. It
records what the implementation already does that supports compliance, and
what is missing.

\paragraph{Authentication.} \emph{Implemented:} hashed passwords, uniform
sign-in errors, rate-limited credential endpoints, and immediate effect of
suspension. \emph{Missing:} revocation of individual sessions, multi-factor
authentication and email verification.

\paragraph{Authorization.} \emph{Implemented:} server-side enforcement of role,
ownership, lifecycle and trust on every request, with scoping inside database
queries and no reliance on client-side checks. Automated tests cover every
role against every scoped read and transition.

\paragraph{Protection of sensitive pickup information.} \emph{Implemented:} a
donor's exact pin, address text and name are shown only to the donor,
organisations able to act on the donation, administrators, and the one courier
who has claimed the run. Other couriers see a roughly one-kilometre area.
Organisation positions are blurred in rankings read by others, and
organisation and courier phone numbers are scoped to readers with a working
relationship. \emph{Missing:} the free-text description, photo and notes are
not scoped, so a donor who types an address into the description would expose
it.

\paragraph{Verification and trust controls.} \emph{Implemented:} an
organisation cannot be ranked, see the open pool or accept food until an
administrator verifies it, and loses verification if it changes its name or
location. \emph{Operational requirement:} the legal and reputational value of
this control depends on administrators performing genuine vetting, which is a
process outside the software.

\paragraph{Data minimisation and user rights.} \emph{Implemented:} registration
asks only for what each role needs; a courier's coordinates are write-only; and
the event ledger records server facts, not free-form surveillance.
\emph{Missing and required before public use:} a privacy notice and consent at
sign-up; a way for users to obtain, correct or delete their data (accounts can
currently only be deactivated); a defined retention period for donations,
photos and events; and hosting and backup arrangements that protect stored
data.

\paragraph{Food safety and liability.} The platform records the storage type,
optional preparation time and deadline so that a recipient can decide, and it
does not certify food. The proposal places inspection with the recipient at
handover~\cite{foodlinkProposal}. The current interface does not yet present
terms of use stating this division of responsibility. That would be needed,
together with review of any food-safety obligations applicable to food
donation, before public operation.

\textbf{Assessment:} the security and privacy architecture is already stronger
than a typical prototype and suitable for a supervised pilot with informed
participants. Public operation additionally requires the privacy, consent,
data-rights, retention and terms-of-use measures listed above.

\section{Feasibility Summary}
\label{sec:feas-summary}

\begin{table}[htbp]
\centering
\small
\begin{tabular}{|L{2.05cm}|L{3.3cm}|L{3.2cm}|L{2.4cm}|L{1.95cm}|}
  \hline
  \textbf{Dimension} & \textbf{Implemented} & \textbf{Feasible, not implemented} & \textbf{Future work} & \textbf{Verdict} \\
  \hline
  Technical & Full workflow, 27 routes, 6 tables, 651 passing tests, CI & PostgreSQL, scheduler, exception handler, token revocation & Multi-process deployment, routing, object storage & Feasible \\
  \hline
  Operational & Four portals covering every core step & Cancel, release and account screens; mobile entry routing & Notifications; broader onboarding & Feasible for pilot \\
  \hline
  Economic & Open-source stack; no paid services & Single-server pilot hosting & Messaging and storage costs & Feasible \\
  \hline
  Schedule & Core system by mid-term & P2 fixes, portal controls, deployment configuration & Pilot, metric collection & Feasible \\
  \hline
  Legal, security, privacy & Scoped access, location protection, verification & Security headers, profile bounds & Privacy notice, consent, data rights, retention, terms & Conditional \\
  \hline
\end{tabular}
\caption{Summary of feasibility by dimension}
\label{tab:feasibility}
\end{table}

Overall, FoodLink is feasible. Its core has moved from proposal to working,
tested software. What remains is predominantly deployment, operational
completeness and legal readiness for real users, rather than open technical
questions.
